Let \(\mathcal I\) be any menu in the stated class.  Simultaneous
identification includes
\(\operatorname{GID}_{\mathcal I}(Q_{0,0})\) and
\(\operatorname{GID}_{\mathcal I}(Q_{1,0})\).  The only-if direction of
\(\texttt{thm:bow-rowspan-iff}\), applied separately to these two queries,
gives
\[
 r_0,r_1\in\operatorname{row}(A_{\mathrm b}).
\]
The treatment-effect span premise is marginal effect completeness in the
binary treatment space.  Therefore
\(\texttt{thm:binary-bow-universal-rank-lower-bound}\) applies and yields
\[
 \operatorname{rank}(A_{\mathrm b})\ge3.
\]
Since each tagged treatment branch contributes one row, every menu in the
class also has at least three tagged treatment branches.

Conversely, \(\texttt{thm:binary-bow-query-separation}\) proves that the
single-instrument menu \(\mathcal I_{\mathrm b}^{\dagger}\) has complete
treatment effect and state spans, has a terminal effect basis, has rank
exactly three, and identifies \(Q_{z^{\circ},y}\) for all
\(z^{\circ},y\in\{0,1\}\) throughout the stated primitive-positive
edge-source class.  Thus attainable ranks are bounded below by three and
include three, proving that the minimum is exactly three.

The marginal-effect-completeness premise is load-bearing.  If it is removed,
take at \(X_{\mathrm b}\) two one-branch instruments, both with effect
\(\mathbf 1\), and with states \(e_0\) and \(e_1\), respectively.  Each
instrument is normalized, but the treatment effects span only
\(\operatorname{span}\{\mathbf 1\}\).  By
\(\texttt{def:bow-probe-matrix}\) and
\(\texttt{def:bow-target-coefficient}\), the treatment matrix has exactly
the two rows \(r_0,r_1\) and rank two.  Paired with the terminal effect
basis, \(\texttt{thm:bow-rowspan-iff}\) identifies all four atomic queries.
This witness shows why the rank-three lower bound is restricted to the
marginally effect-complete class.  The result asserts necessity and explicit
attainment within that class; it does not assert that rank three alone is
sufficient.